三种条件均完成四轮问题。前两轮的六个结果图层逐格通过独立核验，最大长度绝对差约为 $5.1\times10^{-11}$ m，低于预设容差；对应的地图均使用本轮计算图层和正确指标。完整模式用百分比表示暴露率，另外两组用 0--1 分数；核验先统一单位再比较，未把字段名称不同误判为计算错误。

\begin{table}[H]
\centering
\caption{探索性三条件对照。每条件仅一条四轮轨迹；“边界”指水深与封路两问，未做显著性检验。}
\label{tab:gis-comparison}
\small
\begin{tabular}{lrrrrr}
\toprule
条件 & 数值任务 & 边界回答 & 调用数 & 读图数 & 总时长(s)\\
\midrule
依赖追踪 + 图像 &2/2&2/2&27&2&194.7\\
通用工具 + 图像 &2/2&2/2&28&2&200.3\\
依赖追踪，无图像 &2/2&2/2&28&0&185.4\\
\bottomrule
\end{tabular}
\end{table}

上述结果没有显示依赖追踪或图像输入在这批任务上提高成功率。通用条件自行正确重算了新道路集合和新网格，未尝试复用失效结果；完整模式因此没有在自然轨迹中触发旧结果阻断。带图像的两组都调用了截图读取，但没有基于截图再次修改地图。无图像模式也导出了字段、单位与范围正确的地图，并主动说明未进行图面目视检查。不能因为实现了某项机制就认定该机制已经产生增益。

六张地图均有图例、比例尺、网格北、坐标和来源。这里按要素与指标核对地图，未做视觉质量排名。完整结果可在 \code{event/runs/gis-comparison/} 查看。

三组对最大水深均回答现有图层不足、应记为未知而非零；对实际封路长度也明确区分 443.52 m 的平面暴露与缺少通行观测的未知长度。两问没有新增工具调用，说明回答来自已经提供的数据说明和历史工具上下文，并非系统另行获取了水深或交通信息。此处通过只意味着没有超出当前证据，不证明模型能够在未提示限制的开放问题中始终正确拒答。

供应商响应报告的累计 token 用量依次为 512\,381、518\,282 和 456\,904，包含每步重复输入的上下文及缓存 token，不是互不重复的文本量，也不等同于账单费用。时长含接口等待与本地处理，只有单次观测，不能据十余秒差异主张某种模式稳定更快。所有调用均未返回工具错误；这也意味着本组自然任务没有覆盖复杂错误恢复，相关机制仅由单独的故障测试检查。
